\section{Introduction}
\label{sec:intro}

The rapid proliferation of task-specific expert models has made parameter-efficient multi-task synthesis a cornerstone of modern machine learning. Under the paradigm of \emph{model merging} \cite{Wortsman2022Soups, Ilharco2022TaskArithmetic}, multiple neural networks fine-tuned from a shared pre-trained base can be combined directly in parameter space without requiring joint multi-task retraining or accessing the original datasets. This provides a highly scalable and computationally efficient alternative to traditional multi-task training.

Despite its benefits, deploying merged models onto resource-constrained edge devices necessitates compression. \emph{Post-Training Quantization} (PTQ) \cite{Jacob2018Quant, Banner2019PTQ4bit} is the industry standard for reducing memory footprints and accelerating execution speed, translating 32-bit floating-point (FP32) weights into compact, low-precision integers (e.g., 8-bit or 4-bit). However, weight quantization introduces high-frequency rounding noise that corrupts the delicate task-specific vector directions. This corruption is particularly severe in extreme compression regimes (such as 4-bit quantization), where the joint degradation often completely collapses the multi-task capacity of the merged model.

To alleviate this, recent works have proposed optimizing layer-wise merging coefficients at test-time to adapt to quantization noise \cite{Yang2023AdaMerging}. Yet, a fundamental question remains unexplored: \textbf{How does the geometric landscape flatness of the task-specific expert models govern their downstream resilience to post-training quantization and test-time coefficient optimization?}

\begin{figure*}[t]
  \centering
  \begin{subfigure}[b]{0.32\textwidth}
    \includegraphics[width=\textwidth]{flatq_merge_acc_8bit.png}
    \caption{8-bit Quantization Accuracy}
    \label{fig:acc_8bit}
  \end{subfigure}
  \hfill
  \begin{subfigure}[b]{0.32\textwidth}
    \includegraphics[width=\textwidth]{flatq_merge_acc_4bit.png}
    \caption{4-bit Quantization Accuracy}
    \label{fig:acc_4bit}
  \end{subfigure}
  \hfill
  \begin{subfigure}[b]{0.32\textwidth}
    \includegraphics[width=\textwidth]{flatq_merge_curvature_profile.png}
    \caption{Descent-Slope Curvature Profiling}
    \label{fig:curvature}
  \end{subfigure}
  \caption{Overview of our exhaustive empirical findings across 3 independent random seeds. (a) Under low noise (8-bit), multi-task accuracy is stable for SAM radii up to $\rho=0.05$. (b) Under extreme noise (4-bit), we discover a powerful \emph{Flatness-Robustness Synergy}, peaking at $\rho=0.05$ with up to $+7.44\%$ absolute accuracy gains. Both settings degrade rapidly beyond the Over-Perturbation Threshold ($\rho \ge 0.1$). (c) Curvature profiling of the test-time adaptation landscape under varying Gaussian coefficient perturbations ($\sigma$) confirms a stable optimization landscape.}
  \label{fig:main_results}
\end{figure*}

In this paper, we present \textbf{FlatQ-Merge} (Flatness-Aware Quantization-Aware Model Merging), an exhaustive empirical study and framework for analyzing the relationship between expert loss landscape geometry and low-precision multi-task performance. Rather than promoting a single optimization method, our work serves as a systematic comparative study of test-time optimization pathways (direct quantized adaptation vs. unquantized post-hoc adaptation) under different pre-merging expert flatness constraints. Guided by the philosophy of \emph{empirical validation}, we reject minor theoretical conjectures in favor of robust, large-scale, and multi-axial grid sweeps. We pre-train task-specific experts using Sharpness-Aware Minimization (SAM) \cite{Foret2020SAM} across 5 distinct perturbation radii ($\rho \in \{0.0, 0.01, 0.05, 0.1, 0.2\}$). We then merge and quantize these experts to 8-bit and 4-bit precision, comparing 4 different merging methodologies across 3 independent random seeds.

Our extensive empirical evaluation on a Vision Transformer (\texttt{vit\_tiny\_patch16\_224}) \cite{Dosovitskiy2020ViT} across MNIST, FashionMNIST, CIFAR-10, and SVHN yields several highly impactful insights:
\begin{itemize}
  \item \textbf{Precision-Dependent Synergy:} Promoting wider minima during expert pre-training directly improves resilience to quantization noise. Under standard 8-bit quantization, flatness yields negligible gains; however, under extreme 4-bit quantization, pre-training with an optimal SAM radius ($\rho=0.05$) yields a remarkable \textbf{+7.44\%} absolute multi-task accuracy improvement over sharp experts.
  \item \textbf{Flatness Dominates Adaptation:} Merging flat experts with static uniform weights ($\rho=0.05$, NaiveUniform) outperforms sophisticated test-time adaptation on sharp experts ($\rho=0.0$, FlatQ-Merge) by \textbf{+6.03\%} absolute accuracy. This confirms that pre-merging landscape conditioning is more critical than downstream adaptation algorithms.
  \item \textbf{Peak Memory Trade-offs:} While optimizing in FP32 post-hoc slightly outperforms direct quantized optimization under extreme noise, it increases peak RAM by up to 8$\times$. Direct quantized optimization (FlatQ-Merge) keeps weights in 4-bit memory throughout adaptation, maximizing systems efficiency.
  \item \textbf{Over-Perturbation Threshold:} We identify a distinct, non-linear degradation boundary at $\rho \ge 0.1$. Beyond this point, excessive SAM perturbations destabilize pre-training, causing underlearning and low-quality task vectors that collapse multi-task capacity.
  \item \textbf{Curvature Profiling Stability:} Fine-grained coefficient-space sweeps reveal that the test-time adaptation loss landscape is highly stable and robust to perturbations up to $\sigma=0.1$ in 8-bit quantized space, explaining the reliable convergence of gradient updates.
\end{itemize}

By offering overwhelming empirical evidence across different seeds, quantization levels, and tasks, we demonstrate that expert flatness is a primary driver of success in quantized weight spaces, presenting a concrete path forward for deploying robust, compressed, and parameter-fused models.
